\chapter{Метод контекстної оцінки інформаційного ризику}
\label{ch:risk}

% ============================================================
%  РОЗДІЛ 3 -- перший центральний метод середовища виконання політик AURA.
%  Оцінює РИЗИК (дію приймає розділ 4).
%  Психологічні поняття -- лише як формалізовані сигнали.
% ============================================================

\section{Повідомлення як подія в контексті, а не ізольований текст}
\label{sec:event-not-text}
% -- Формалізація події: (відправник, контакт, час, історія, домен).

\section{Контактна пам'ять}
\label{sec:contact-memory}
% -- Стан відправника/контакту; довіра, попередні сигнали.

\section{Пам'ять розмови}
\label{sec:conversation-memory}
% -- Локальний стан діалогу; перехід тем; контекстні маркери.

\section{Часові закономірності, повторюваність та ескалація}
\label{sec:temporal}
% -- Частота, час доби, сплесковість, наростання тиску.

\section{Латентні стани маніпуляції (операціоналізація)}
\label{sec:latent-states}
% -- грумінг, примус і маніпуляція як формалізовані сигнали, НЕ діагноз.

\section{Доменні модулі ризику}
\label{sec:domain-modules}
% -- дитячий, військовий та інші предметно-орієнтовані режими; конфігурованість.

\section{Формування структурованого сигналу ризику}
\label{sec:risk-signal}
% -- Вихід методу: структурований сигнал ризику для розділу 4.
Результатом методу є структурований сигнал ризику, що інтегрує контактну пам'ять, пам'ять розмови, часову динаміку та доменний контекст і слугує входом для прийняття рішень з урахуванням політик (розділ~4). \emph{[Розкрити формалізацію та структуру сигналу ризику.]}

\rozdilconcl{3}
\emph{[Висновки до розділу 3.]}
